\documentclass[11pt,a4paper]{article}

\usepackage[margin=1in]{geometry}
\usepackage{lmodern}
\usepackage[T1]{fontenc}
\usepackage{microtype}
\usepackage{booktabs}
\usepackage{parskip}
\usepackage[hidelinks]{hyperref}
\usepackage{titling}

\setlength{\droptitle}{-2.5em}

\title{\Large FinXplain: Request for Expert Review of an Explainable AI Prototype \\[2pt]
       \large for Philippine Fintech Savings-Product Recommendations}
\author{Ashton Stephonie T. Matias \\
        Bachelor of Science in Computer Science \\
        University of the Philippines Los Ba\~{n}os \\
        \texttt{atmatias@up.edu.ph}}
\date{\today}

\begin{document}

\maketitle
\thispagestyle{empty}

\section*{Purpose of this brief}

This document requests your expert review of \textbf{FinXplain}, an undergraduate
research prototype that recommends Philippine savings and investment products
to retail users together with a plain-language explanation for each
recommendation. I am writing to ask whether the product--profile pairings the
system produces are consistent with sound personal-finance practice as you
would advise a Filipino client today. Your feedback will be cited in the final
manuscript and acknowledged in the defense.

The prototype is a \emph{research artifact} only. It is not deployed to the
public, does not connect to any provider API, and is not intended to issue
binding financial advice; the user-facing interface displays this disclaimer
prominently.

\section*{What the system does, in one paragraph}

FinXplain takes 16 inputs from the user (age, monthly income, monthly
expenses, existing savings, employment status, educational attainment, number
of dependents, location type, digital savviness, bank-account ownership,
e-wallet ownership, primary e-wallet, OFW remittance receipt, savings goal,
risk tolerance, and investment horizon) and recommends one of eight savings
or investment products from three providers (GCash, Maya, BPI). For every
recommendation it returns: (i)~the predicted product and the top-3
alternatives with probabilities, (ii)~a ranked top-5 list of factors that
\emph{supported} or \emph{opposed} the recommendation, (iii)~plain-language
sentences explaining each factor in everyday terms, and (iv)~up to three
\emph{counterfactual} suggestions of the form ``if you increased your
existing savings to PHP~50{,}000, BPI Plan Ahead Time Deposit would be
recommended instead.''

The recommendation engine is an Explainable Boosting Machine (EBM)---a
glass-box model whose contributions per feature can be read off directly,
without post-hoc approximation. An XGBoost benchmark is used only for
offline performance comparison.

\section*{Product catalogue under review}

The eight products covered by the prototype are listed in
Table~\ref{tab:products}. Interest rates and minimum balances are point-in-time
values verified in April 2026 against the providers' official documentation.

\begin{table}[h]
    \centering
    \small
    \caption{Products covered by the FinXplain prototype.}
    \label{tab:products}
    \begin{tabular}{lllll}
        \toprule
        \textbf{Product} & \textbf{Provider} & \textbf{Type} & \textbf{Posted Rate} & \textbf{Min. Balance} \\
        \midrule
        GSave (CIMB)             & GCash & Digital savings        & 2.6\% p.a.          & PHP 0       \\
        GInvest (UITFs)          & GCash & Investment (UITF)      & Market-linked       & PHP 50      \\
        Maya Savings             & Maya  & Digital savings        & 3.0\% base, up to 15\% & PHP 0    \\
        Maya Personal Goals      & Maya  & Goal-based savings     & 4--8\% tiered       & PHP 0       \\
        Maya Time Deposit Plus   & Maya  & Digital time deposit   & 3.5--6.0\% p.a.     & PHP 5{,}000 \\
        BPI Regular Savings      & BPI   & Traditional savings    & 0.0625\% p.a.       & PHP 3{,}000 \\
        BPI \#SaveUp             & BPI   & Digital auto-save      & 0.0925\% p.a.       & PHP 1       \\
        BPI Plan Ahead TD        & BPI   & 5-year time deposit    & Branch-quoted fixed & PHP 50{,}000\\
        \bottomrule
    \end{tabular}
\end{table}

\section*{How the system decides}

Because Philippine fintech transaction data is not publicly available, I
generated a synthetic dataset of 5{,}000 profiles whose marginal and
conditional distributions are parameterized from PSA FIES 2023, PSA LFS
2023, the PSA 2020 Census, and the BSP 2021 Financial Inclusion Survey. I
then assigned a product label to every profile through a hand-coded,
multi-factor additive scoring system grounded in each product's official
eligibility criteria and in established personal-finance principles. The
EBM is trained on these (profile, product) pairs.

A condensed version of the labeling rules is shown in
Table~\ref{tab:rules}. Note that these rules are deliberately a coarse
caricature of the suitability assessment a real advisor would perform; the
specific question I am asking you is whether the resulting product--profile
pairings are nonetheless defensible, and where you would push back.

\begin{table}[h]
    \centering
    \footnotesize
    \caption{Condensed rule logic for product assignment (R1--R10).}
    \label{tab:rules}
    \begin{tabular}{p{2.4cm}p{12cm}}
        \toprule
        \textbf{Product} & \textbf{Profile (informal)} \\
        \midrule
        GSave (GCash)            & Low existing savings ($<$PHP 5{,}000); digitally engaged; primary e-wallet GCash. \\
        GInvest (GCash)          & Aggressive risk tolerance; long horizon; college-or-graduate education boosts further. \\
        Maya Savings             & Low-to-moderate income; primary e-wallet Maya or Both; digitally engaged. \\
        Maya Personal Goals      & Education / Home-Property / Travel goals; medium horizon; remittance recipient boosts. \\
        Maya Time Deposit Plus   & Moderate income (PHP 30k--60k); medium horizon; non-aggressive risk; college-plus boosts. \\
        BPI Regular Savings      & Has bank account; conservative; lower digital savviness; e-wallet ``Neither'' boosts. \\
        BPI \#SaveUp             & Has bank account; digitally engaged; non-retirement goal; remittance-recipient boosts. \\
        BPI Plan Ahead TD        & Existing savings $\geq$ PHP 50{,}000; conservative; long horizon; college-plus boosts. \\
        \bottomrule
    \end{tabular}
\end{table}


\section*{What I am asking}

The defense panel for this study recommended that I consult a practising
financial advisor on the substantive content of the recommendations,
separately from the machine-learning evaluation. Specifically, I would value
your view on the following five questions:

\begin{enumerate}
    \item \textbf{Suitability.} Looking at the rule logic in
    Table~\ref{tab:rules}, are these product-to-profile assignments
    consistent with what you would actually advise a Filipino client in
    each profile category? Where would you disagree?
    \item \textbf{Missing factors.} What additional inputs (e.g.\
    insurance coverage, outstanding debt, irregularity of income) do you
    consider essential for a responsible savings or investment
    recommendation, and which of these would you prioritize adding to a
    next iteration?
    \item \textbf{Counterfactual responsibility.} The system suggests
    actionable changes such as ``increase your existing savings to PHP
    X'' to unlock a different product. Are there counterfactual
    suggestions that you consider inappropriate to surface to retail
    users (e.g.\ those that imply unrealistic short-term income changes)?
\end{enumerate}

\section*{How your feedback will help}

Even short, candid reactions would be enormously helpful---I am not looking
for a formal report. \textbf{A few bullet points per question, or even a
quick voice note, is more than enough}; please feel free to skip any
question that doesn't apply to your practice, and if something else jumps
out that I didn't think to ask, I would love to hear about that too.

Whatever you share will go toward (i)~tightening the rule logic before
final manuscript submission, (ii)~a short \emph{Expert Review} subsection
in the Discussion chapter that summarizes your input and the changes I
made in response (or, where I kept something as-is, why), and (iii)~the
disclosure language shown to users on the prototype's results page. With
your permission, I would be honoured to acknowledge you by name and
affiliation in the final paper.

If it is easier than typing, I am happy to walk you through the prototype
live in a 20--30 minute call or visit at any time that fits your schedule.
Otherwise, a reply by email---bullet points, voice notes, or even margin
comments on this PDF---is perfect. There is no rigid deadline on my side;
I will adapt to whatever pace works for you. Thank you so much for
considering this---your time and expertise would genuinely shape the final
version of this work.

\vspace{1.5em}

\noindent Sincerely,

\vspace{2em}

\noindent Ashton Stephonie T. Matias \\
\noindent BS Computer Science, UP Los Ba\~{n}os

\end{document}
